\chapter{Expression complexity and resource semantics}
A system stores four experts but routes each token to one. Has it reduced
inference cost by a factor of four? Only if we know what is counted.
The router still runs; weights still occupy storage; padded execution may
evaluate empty slots; moving representations may dominate arithmetic.
A sparse logical graph is not a measured speedup.

Chapter 15 gave executions an exact transition semantics. Here we attach
explicit resource models to computations. We use integer ledgers and a tiny
dependency graph, not GPU timings. The purpose is to make a cost claim
falsifiable: a reader should be able to identify both the counted work and the
work omitted by the model.

\section{A resource is a vector with units, not one adjective}
\index{resource semantics}\index{cost model}
For an execution trace $\tau$, define separate components
\begin{equation}
 C(\tau)=(P,\ M,\ R,\ D,\ H,\ J).
\end{equation}
Here $P$ counts stored scalar parameters, $M$ multiply-accumulate operations
(MACs), $R$ routing decisions, $D$ communicated bytes under a stated interface,
$H$ peak live buffer bytes, and $J$ retained events.
These quantities cannot be added without conversion factors and a purpose.
A weighted objective $w^\top C$ introduces policy weights; it is not a universal
physical cost.

\Plate{01-ledger}{Resource ledgers have different units and aggregation rules.
Shared ownership can deduplicate stored parameters while repeated calls still
perform arithmetic. Peak memory is a maximum over live sets, not a sum over
the entire trace.}{fig:ledger}

For additive events, $M(\tau_1\tau_2)=M(\tau_1)+M(\tau_2)$.
Peak storage needs the intermediate live state, so it cannot generally be
computed from two isolated peaks alone. Parameter storage counts unique owners,
not appearances in an instruction list. These distinctions prevent a repeated
shared module from being mistaken for either free work or duplicated weights.

Every ledger needs a boundary. Does memory include optimizer state, activations,
input storage, allocator fragmentation and kernel scratch space? Does communication
count one-way payload, both directions, padding or protocol overhead?
We will answer narrowly before computing anything.

\section{Derive an expert ledger from tensor dimensions}
\index{mixture of experts}\index{multiply-accumulate}
Consider $B$ tokens of width $d$, $E$ experts, hidden width $h$, and top-$k$
routing without repeated expert IDs per token. Each expert has bias-free
matrices $W_1\in\mathbb R^{d\times h}$ and
$W_2\in\mathbb R^{h\times d}$. The router has $W_R\in\mathbb R^{d\times E}$.
Thus stored parameter count is
\begin{equation}
 P=2Edh+dE.
 \label{eq:params}
\end{equation}
We count each multiply-accumulate as one MAC, including accumulation initialized
at zero. Applying one expert to one token costs $2dh$ MACs, while router
logits cost $BdE$ MACs. Selection, activation, normalization and aggregation
are outside this MAC count and must not be assumed free in elapsed time.

Without drops or padding, useful expert work is
\begin{equation}
 M_{\rm useful}=2Bkdh.
 \label{eq:useful}
\end{equation}
If only $U$ unique experts are called, their expert parameter count is $2Udh$.
This is an active-owner count, not the number of parameters stored in the model
and not a guarantee those bytes are fetched just once.

Suppose each expert is allocated $C$ token slots. With dense padded kernels
over every expert's slots, executed expert work is
\begin{equation}
 M_{\rm padded}=2ECdh.
 \label{eq:padded}
\end{equation}
Let $n_e$ count requested assignments to expert $e$. Accepted work uses
$\min(n_e,C)$ assignments, and dropped assignments total
$\sum_e\max(0,n_e-C)$. The reference drops overflow in input order; different
admission policies change which tokens lose work.

\Plate{02-routing}{Four tokens request experts 0,0,0,1. A padded grid reserves
slots for every expert, including unused ones. Capacity two drops one request
to expert 0; it is not merely an implementation detail if output semantics
depend on that assignment.}{fig:routing}
\Listing{routing_cost}

For $B=4,E=4,d=8,h=16,k=1$, the model stores 1056 scalars:
1024 expert and 32 router parameters. Router work is 128 MACs.
The two used experts own 512 scalars. Generated expert-work counts are
\begin{center}\input{\Generated/routing.tex}\end{center}
The padded $C=4$ case executes four times the useful expert MACs. The $C=2$
case drops one assignment yet still executes more expert MACs than the ragged
case. Fewer accepted assignments can coexist with more executed work.

For a float32 dispatch interface that sends each accepted width-$d$ activation
and receives one width-$d$ result, payload bytes are
\begin{equation}
 D=2\cdot4d\sum_e\min(n_e,C),
 \label{eq:bytes}
\end{equation}
using $n_e$ when capacity is unbounded. The ragged example gives 256 bytes.
This excludes parameter transfer, routing metadata, collective overhead and
padding on the wire. Local assignments may need no inter-device traffic at all;
the formula describes a declared all-remote payload model, not measured network load.

\section{Liveness turns a graph into a memory calculation}
\index{liveness}\index{scheduling}
Consider dependencies $a=f(x)$, $b=g(x)$, $c=u(a)$, $d=v(b)$,
$y=w(c,d)$. Let $x,a,b$ each occupy 100 abstract bytes and $c,d,y$ each
occupy 10. These are invented allocation sizes, not tensor measurements.
Inputs count as live from the beginning. Outputs remain live at the end.
Each result is allocated before last-use operands are released.

A legal schedule must execute a node after its dependencies. Two legal choices are
\[
 S_1=(a,b,c,d,y),\qquad S_2=(a,c,b,d,y).
\]
For $S_1$, allocating $b$ temporarily leaves $x,a,b$ live: 300 bytes.
For $S_2$, computing $c$ permits release of $a$ before allocating $b$:
the peak is 210 bytes. Both schedules execute every operation exactly once.

\Plate{03-liveness}{The same dependency graph admits different legal buffer
lifetimes. The interleaved schedule shortens the lifetime of a large intermediate.
Allocation occurs before last-use release, so the peak includes both old
operands and a newly allocated result.}{fig:liveness}

\Listing{peak_live}

The companion's \texttt{validate\_layout} checks declarations and schedule
membership before execution. The algorithm initializes remaining-use counts
from all dependency edges.
After a node runs, decrement each operand's count and release non-output operands
whose count reaches zero. Multiple uses of the same operand in one node count
as separate uses but do not allocate duplicate buffers.
A node with no future users is released immediately unless declared an output.
Unused initial inputs remain retained in this narrow contract.

To validate the implementation independently, tests enumerate all legal
topological schedules and compute interval overlap from producer and last-user
indices. Both methods must agree; enumerated peaks range from 210 to 300 bytes.
An invalid schedule is rejected rather than assigned an artificially small peak.

This analysis omits aliasing, in-place updates, recomputation, asynchronous
execution, allocator reuse and operator workspace. An actual runtime may save
or spend more memory. In-place reuse is legal only when alias and effect
conditions hold; overwriting an operand before its final consumer changes the
computation. Checkpointing can lower retained activation memory by recomputing
values, changing both the arithmetic ledger and the schedule.

\section{Amortized cost explains totals without hiding spikes}
\index{amortized analysis}\index{potential function}
A trace recorder appends events to a dynamic buffer. Start with size $n=0$
and capacity $c=0$. If full, allocate capacity $\max(1,2c)$ and copy all $n$
old events, then write the new event. Count one unit per event copy or append;
allocation overhead and variable-size event payload are excluded.

An append normally costs one unit. When the old buffer is full at size $n>0$,
it costs $n+1$. The next append can therefore be expensive even if average
cost remains bounded.
Define a nonnegative potential over reachable buffer states:
\begin{equation}
 \Phi(n,c)=2n-c.
\end{equation}
At empty state it is zero. After the first allocation it is one; after later
doublings it is two; between doublings it rises by two per append.
The amortized cost is
\begin{equation}
 \widehat a_i=a_i+\Phi_i-\Phi_{i-1}.
 \label{eq:amort}
\end{equation}
For a non-resizing append, $\widehat a_i=1+2=3$.
For a resize with old $c=n>0$, old potential is $n$ and new potential is 2,
so $\widehat a_i=(n+1)+2-n=3$.
The first append has amortized cost 2.

\Plate{04-potential}{A doubling buffer trades occasional copy bursts for bounded
aggregate work. The potential records prepaid work in the proof; it is not
stored energy or a measured device counter.}{fig:potential}
\Listing{trace_buffer}
\begin{center}\input{\Generated/buffer.tex}\end{center}

Summing Equation~\ref{eq:amort} telescopes:
\begin{equation}
 \sum_{i=1}^{N}a_i
 =\sum_{i=1}^{N}\widehat a_i+\Phi_0-\Phi_N
 \leq 3N.
 \label{eq:aggregate}
\end{equation}
Nonnegative final potential is essential to the bound. Tests check every prefix
through 128 events, not only powers of two.

Modern cost-aware verification work explicitly relates a program's behavior
to its resource accounting \cite{potential}. Our dynamic-buffer example is
an original small application of the potential method, not an implementation
of an automatic resource-analysis tool.
An amortized bound is neither a probabilistic average over workloads nor a
worst-case per-event latency bound. A latency-sensitive recorder may preallocate,
chunk storage or stream events; each alternative changes the contract.

The Chapter 15 trace contains complete environments, not fixed-size events.
If each event stores $r$ registers of up to $b$ bits, its payload scales as
$O(rb)$ before serialization overhead. The buffer's copy-count model cannot
be reported directly as total bytes or constant-time logging.

\section{A performance envelope is a hypothesis to measure}
\index{arithmetic intensity}\index{roofline}
Let $F$ be counted floating-point operations, $D$ actual bytes transferred
across one specified memory boundary, $P_{\max}$ an attainable compute ceiling
in operations per second, and $\beta$ an attainable bandwidth in bytes per
second. Ignoring overlap overhead and other bottlenecks yields the necessary bound
\begin{equation}
 t\geq\max\left(\frac{F}{P_{\max}},\frac D\beta\right).
 \label{eq:time}
\end{equation}
For $D>0$, intensity $I=F/D$ gives
\begin{equation}
 \frac Ft\leq\min(P_{\max},\beta I).
 \label{eq:roof}
\end{equation}
The crossover is $I^*=P_{\max}/\beta$. Below it, bandwidth limits this
two-resource envelope; above it, compute does. Counting a MAC as two operations
is a reporting convention and must be stated consistently.

\Plate{05-roofline}{An illustrative envelope with compute ceiling 100 Gop/s and
bandwidth 10 GB/s has crossover 10 operations/byte. These assigned values are
not a hardware benchmark or a prediction for either Evolutor model direction.}{fig:roofline}

Equation~\ref{eq:time} is a lower bound on time, not a promise that either resource
can be saturated. Launch overhead, dependency chains, synchronization,
small matrices and communication can make observed time much larger.
Do not insert the dispatch-only byte count from Equation~\ref{eq:bytes}
as total device-memory traffic without evidence: it crosses a different boundary.

Current PyTorch profiler documentation provides execution activity and memory
recording, and limits its FLOP estimates to supported operations \cite{profiler}.
Use profiling to test a specified hypothesis after validating output parity.
Declare warm-up, input shapes, dtype, device, synchronization, compilation state,
sample count and uncertainty. Do not compare an eager reference with a compiled
candidate and attribute every difference to the mathematical architecture.

A useful experimental sequence is: establish semantics; predict the ledger;
measure actual operators and traffic; identify the gap; change one implementation
choice; rerun correctness and measurements. This turns a resource estimate into
an engineering question rather than an industrial performance claim.

\section{Exercises with worked answers}
\begin{enumerate}
\item \textbf{Count.} Recompute stored parameters for $E=4,d=8,h=16$.
\emph{Answer:} $2Edh+dE=1024+32=1056$.
\item \textbf{Separate.} Why are 512 active expert scalars not the model size?
\emph{Answer:} Only two experts are called, but all four plus router remain stored.
\item \textbf{Derive.} What is ragged expert work for four top-1 tokens?
\emph{Answer:} $2(4)(1)(8)(16)=1024$ MACs, excluding router and non-MAC work.
\item \textbf{Padding.} What does capacity four execute?
\emph{Answer:} $2(4)(4)(8)(16)=4096$ expert MACs under all-expert dense padding.
\item \textbf{Overflow.} How many assignments are dropped with loads $(3,1,0,0)$
and capacity two? \emph{Answer:} One; capacity does not change the fact that its
omitted contribution can change the output.
\item \textbf{Units.} Verify the ragged dispatch payload.
\emph{Answer:} Four accepted calls times two directions times eight float32 scalars
times four bytes gives 256 bytes.
\item \textbf{Schedule.} Why is $a,c,b,d,y$ better for the example's peak?
\emph{Answer:} It releases 100-byte $a$ after producing 10-byte $c$ before
allocating $b$, reducing the peak from 300 to 210.
\item \textbf{Oracle.} Independently check liveness.
\emph{Rubric:} Enumerate legal orders, assign allocation/last-use intervals,
include initial inputs and final outputs, and compare peaks at allocation boundaries.
\item \textbf{Potential.} Bound total append/copy work.
\emph{Answer:} Telescoping and $\Phi_N\geq0$ give at most $3N$ counted units.
\item \textbf{Latency.} Does amortized cost three bound every append by three?
\emph{Answer:} No; appending to a full capacity-eight buffer copies eight old events
and writes one, costing nine.
\item \textbf{Envelope.} Find crossover for 100 Gop/s and 10 GB/s.
\emph{Answer:} Ten operations/byte. Consistent decimal giga units cancel.
\item \textbf{Experiment.} Design an honest sparse-versus-dense comparison.
\emph{Rubric:} Match output contract, token cohort, precision and hardware; count
router, padding, transfers and compilation; validate parity; report repeated
timings and memory, including failures and dropped assignments.
\end{enumerate}
\section{Terms and next scope}
A \Term{cost model} chooses counted events and units.
The \Term{liveness} of a value depends on its future consumers.
An \Term{amortized cost} bounds aggregate sequences using a proof, not timing.
The \Term{arithmetic intensity} is operations per byte across a specified boundary.
Using \Term{padding} can turn sparse useful work into denser executed work.
The README reproduces all ledgers. Later engine chapters can replace these
assumptions with measurements; they must preserve the semantic contracts first.
